\section{Cyclotomic towers and a common interval for the cohomological bounds}
\label{sec:ct-intervals}

We reconstruct the part of the cohomological argument\cite{alpo-uniform-source} that treats the
prime $2$ and a fixed finite set of odd primes. The algebraic estimates
below allow finite torsion in the modules. The only deep classical
input in the uniform generator bound is the following precise form of
the Ferrero--Washington theorem: if $K/\Q$ is finite abelian, the
$p$-ranks of the ideal class groups of the fields
$K(\mu_{p^n})$ are bounded as $n$ varies. Here the $p$-rank of a finite
abelian group $A$ means $\dim_{\mathbb F_p} A/pA$.
This is the usual consequence of the vanishing of the Iwasawa
$\mu$-invariant and the structure theorem for the unramified Iwasawa
module; we use it as a classical input\cite{ct-fw,ct-sharifi}.
The final interval theorem is first stated with an explicit coarse
cohomological hypothesis. Theorem~\ref{kt-fourteen} below proves
that hypothesis, using precisely stated classical $K$-theory and
duality theorems.

\subsection{The local conditions and cyclotomic descent}

Fix a prime $p$ and a number field $K$ containing $\mu_p$ if $p$ is
odd, and $\mu_4$ if $p=2$. Put
$K_\infty=K(\mu_{p^\infty})$ and
$\Gamma=\operatorname{Gal}(K_\infty/K)$. The cyclotomic character
$\varepsilon_p$ identifies $\Gamma$ with $1+p^c\Z_p$, where
$c\ge1$, and $c\ge2$ for $p=2$. Equivalently, $c$ is the greatest
integer for which $\mu_{p^c}\subset K$. Choose a topological generator
$\gamma$ and put $u=\varepsilon_p(\gamma)$; thus
$v_p(u-1)=c$. Let $K_m$ be the subfield with
$[K_m:K]=p^m$.

For $e\ge1$, let $A_e=\Q_p/\Z_p(-e)$: its underlying group is
$\Q_p/\Z_p$ and $g$ acts by multiplication by
$\varepsilon_p(g)^{-e}$. Define
\[
 \mathcal S_p(F,e)=
 \ker\left(H^1(G_{F,S},A_e)\longrightarrow
                  \prod_{v\mid p}H^1(F_v,A_e)\right),
 \qquad S=\{v:v\mid p\}\cup\{v:v\mid\infty\}.
\]
Here $G_{F,S}$ is the Galois group of the maximal extension unramified
outside $S$, and $H^1$ is continuous group cohomology. In terms of the
absolute Galois group, this imposes zero restriction to inertia away
from $p$ and zero restriction to the \emph{full decomposition group}
at each prime above $p$. Define
\[
 \lambda_p(F,e)=\length_{\Z_p}\mathcal S_p(F,e)
       \quad\hbox{in }\mathbb N\cup\{\infty\}.
\]
When it is finite, $\#\mathcal S_p(F,e)=p^{\lambda_p(F,e)}$.

Let $D_p(K_\infty)$ be the discrete group of continuous homomorphisms
$G_{K_\infty}\to\Q_p/\Z_p$ that vanish on every finite inertia
group and on every decomposition group above $p$. Its Pontryagin
dual is the compact $\Z_p$-module
\[
 X_p(K_\infty)=
 \operatorname{Hom}_{\Z_p}(D_p(K_\infty),\Q_p/\Z_p).
\]
The topology on this dual is the topology of pointwise convergence
on finite subgroups. The action on $D_p$ is
$(\gamma\varphi)(h)=\varphi(\gamma^{-1}h\gamma)$; the action on
$X_p$ is the contragredient action. We write $T$ for the action of
$\gamma$ on $X_p$.

\begin{proposition}[Descent with the full local conditions]
\label{ct-descent}
For every $e\ge1$ and $m\ge0$, there is a canonical isomorphism
\[
 \mathcal S_p(K_m,e)^\vee\simeq
 X_p(K_\infty)\big/
       \bigl((u^eT)^{p^m}-1\bigr)X_p(K_\infty).
\]
In particular, the two sides have the same $\Z_p$-length, with
infinite length permitted.
\end{proposition}

\begin{proof}
The group $G_{K_\infty}$ acts trivially on $A_e$.
For a procyclic pro-$p$ group generated by $\tau$, continuous
cohomology on discrete $p$-primary modules is computed in degrees
$0$ and $1$ by $B\xrightarrow{\tau-1}B$, and vanishes above
degree $1$. Indeed, its completed group ring is $\Z_p[[t]]$,
with $t=\tau-1$, and
$0\to\Z_p[[t]]\xrightarrow{t}\Z_p[[t]]\to\Z_p\to0$
is the continuous free resolution of the trivial module.
For $B=A_e$, multiplication by
$u^{-ep^m}-1$ is surjective on $\Q_p/\Z_p$, since this is a
nonzero scalar and that group is divisible. The scalar is nonzero
because $u$ has infinite order and $e\ne0$. Inflation--restriction
therefore gives
\[
 H^1(K_m,A_e)\simeq
 H^1(K_\infty,A_e)^{\Gamma^{p^m}}.
\]

At a prime not above $p$, the cyclotomic extension is unramified,
so the inertia group is unchanged on passing to $K_\infty$.
At a prime above $p$, the image of the decomposition group in
$\Gamma$ is open and nontrivial: a finite extension of $\Q_p$
contains only finitely many $p$-power roots of unity. The same
divisibility argument makes the local inflation kernel zero.
Thus a global class has zero localization there if and only if its
restriction to the decomposition group over $K_\infty$ is zero.
The displayed global isomorphism consequently identifies the strict
Selmer group with the invariants of $D_p(K_\infty)(-e)$.

Pontryagin duality is exact here, since $\Q_p/\Z_p$ is injective
as a $\Z_p$-module. The dual of invariants is the corresponding
group of coinvariants. Dualizing the twist $-e$ replaces it by $e$;
the action on the dual is therefore $u^eT$. This proves the formula.
It also proves the asserted equality of lengths, by dualizing finite
composition series; either side has finite length precisely when
the other does.
\end{proof}

\subsection{A uniform bound as the quadratic fields vary}

We first give the elementary group argument that controls the cost
of a quadratic extension. If $V$ is profinite and $S\subset V$ is
stable under conjugation, write
\[
 C(V,S)=\{\varphi:V\to\mathbb F_2:
       \varphi\text{ is a continuous homomorphism},\ 
       \varphi|_S=0\}.
\]

\begin{lemma}[A quadratic extension]
\label{ct-index-two}
Let $U$ be an open subgroup of index $2$ in $V$, and suppose
$S\subset U$ is stable under conjugation by $V$. If
$\dim C(V,S)=r<\infty$, then
\[
                 \dim C(U,S)\le 2r.
\]
\end{lemma}

\begin{proof}
Choose $\sigma\in V\setminus U$. Its action $a$ on
$W=C(U,S)$ is an involution, because $\sigma^2\in U$ and inner
automorphisms act trivially on homomorphisms to an abelian group.
In characteristic $2$, $(a-1)^2=0$, so
$\operatorname{im}(a-1)\subset W^a$. It remains to show
$\dim W^a\le r$.

For $\varphi\in W^a$, the only obstruction to extending
$\varphi$ to $V$, with value $0$ at $\sigma$, is
$\varphi(\sigma^2)$. To check this directly, write each element
of $V$ uniquely as $h\sigma^i$ with $h\in U$ and $i\in\{0,1\}$,
and set $\widetilde\varphi(h\sigma^i)=\varphi(h)$.
The multiplication law gives a homomorphism exactly when
$\varphi$ is invariant and $\varphi(\sigma^2)=0$.
Continuity follows since $U$ is open. It still vanishes on $S$.
The restriction map from $C(V,S)$ to $W^a$ has a one-dimensional
kernel, the character of $V/U$, and its image has codimension at
most one in $W^a$. Hence $\dim W^a\le(r-1)+1=r$.
Finally, $W/W^a$ injects into $W^a$ via $a-1$, proving both
finiteness of $W$ and the asserted bound.
\end{proof}

\begin{lemma}[Descent of unramified characters]
\label{ct-character-descent}
Let $K_\infty/K$ be a cyclotomic $\Z_p$-extension.
Every continuous homomorphism
$G_{K_\infty}\to\mathbb F_p$ unramified at all finite primes
is the restriction of such a homomorphism on $G_{K_n}$ for some
$n$. A finite family can be descended to a common level.
\end{lemma}

\begin{proof}
Put $G=G_K$ and $H=G_{K_\infty}$. Choose an open normal subgroup
$N$ of $G$ such that $N\cap H$ lies in the kernel of the given
homomorphism $\varphi$. The rule $hn\mapsto\varphi(h)$ extends
it to $HN$: the ambiguity and the commutators $[n,h]$ lie in
$N\cap H$, where $\varphi$ vanishes. Since the decreasing compact subgroups $G_{K_n}$
have intersection $H$, some $G_{K_n}$ lies in $HN$.

At primes away from $p$, inertia lies in $H$, so the extension
already vanishes on inertia. The union of the conjugates, inside
$G_{K_n}$, of its finitely many inertia groups above $p$ is compact:
it is a finite union of continuous images of products of compact
groups. If the extension failed to vanish on that union at every
deeper level, fix this compact union $C$ at the initial level and
consider the decreasing compact sets
$C\cap G_{K_m}\cap\widetilde\varphi^{-1}(\mathbb F_p\setminus\{0\})$.
If all were nonempty, they
would have nonempty intersection. That would give an element of
$H$ in an inertia group on which $\varphi$ is nonzero, a
contradiction. A deeper level is therefore unramified also at $p$.
Taking the largest of finitely many levels gives the last assertion.
\end{proof}

\begin{proposition}[Uniform generator bound]
\label{ct-uniform-generators}
Let $M/\Q$ be finite abelian, of degree $d_M$ and conductor $f$.
There is an integer $g$, depending only on $M$, with the following
property. If $q,d$ are primes coprime to $2f$ satisfying
\[
                    q\equiv3\pmod8,
       \qquad       d\equiv7\pmod {16},
\]
and $L=M(i,\sqrt{-q},\sqrt{-d})$, then
$X_2(L_\infty)$ is generated over $\Z_2$ by at most
\[
                         R=4(g+6d_M)
\]
elements. This bound is independent of $q$ and $d$.
For any fixed abelian field $K$ containing the requisite roots of
unity, $X_p(K_\infty)$ is finitely generated over $\Z_p$ as well.
\end{proposition}

\begin{proof}
Apply Ferrero--Washington to the fixed field $M(i)$. By global
class field theory, the dimension of its unramified
$\mathbb F_2$-valued characters at each finite cyclotomic level
equals the $2$-rank of the ideal class group. Choose a uniform
bound $g$. Lemma~\ref{ct-character-descent}, applied to any finite
family, bounds the space of unramified characters of
$M(\mu_{2^\infty})$ by $g$.

The number of primes above an odd prime $a$ in
$\Q(\mu_{2^\infty})$ is the index of the closure of $\langle a\rangle$
in $\Z_2^\times$. For $a\equiv3\pmod8$ this index is $2$;
for $a\equiv7\pmod {16}$ it is $4$. Here is a direct verification.
If $a\equiv3\pmod4$ and $s=v_2(a^2-1)$, then
\[
 v_2(a^{2^k}-1)=s+k-1\quad(k\ge1).
\]
This follows by induction, since each subsequent factor
$a^{2^k}+1$ has valuation $1$. For $n\ge s$, the order of $a$
modulo $2^n$ is $2^{n-s+1}$ and the index is $2^{s-2}$.
The two congruences give $s=3$ and $s=4$, respectively.
Adjoining $M$ increases a number of primes by at most $d_M$.
There are consequently at most $2d_M$ primes above $q$ and
$4d_M$ above $d$ in $M(\mu_{2^\infty})$.

Allow ramification at those primes. The space of characters can
increase in dimension by at most one for each such prime. Indeed,
wild inertia has pro-odd order and every homomorphism from its
tame quotient to $\mathbb F_2$ is determined by one value.
Evaluation on one tame generator at each prime has the unramified
character space as its kernel. Conjugate primes give the same
condition. Thus the dimension after allowing $q,d$ is at most
$g+6d_M$.

Adjoin $\sqrt{-q}$ and then $\sqrt{-d}$. Each extension has
degree at most $2$ and is unramified outside the allowed primes:
the quadratic discriminants are $-q$ and $-d$.
Apply Lemma~\ref{ct-index-two} at every nontrivial step, taking
$S$ to be the union of the inertia groups at the other primes.
The dimension is multiplied by at most $4$. Imposing
unramifiedness again at $q,d$, and complete splitting above $2$,
can only decrease it. Therefore
\[
             \dim_{\mathbb F_2}D_2(L_\infty)[2]\le R.
\]

Duality identifies the dual of this last group with
$X_2(L_\infty)/2X_2(L_\infty)$. A compact pro-$p$ abelian group
$X$ with $\dim X/pX\le R$ is generated by $R$ elements as a
$\Z_p$-module. To see this without a freeness assumption, choose
lifts of a basis and let $Y$ be the image of $\Z_p^R$.
It is compact, hence closed. The quotient $X/Y$ satisfies
$X/Y=p(X/Y)$, so every finite $p$-group quotient is zero.
A profinite group is separated by its finite quotients, giving
$X/Y=0$. This proves the uniform assertion. The same descent
and duality argument, with Ferrero--Washington for a fixed $K$,
proves the last assertion.
\end{proof}

\begin{remark}
The congruences control the numbers of primes in the infinite
cyclotomic tower. Replacing them by mere oddness would remove this
particular uniform bound: for $a\equiv1\pmod {2^k}$ the closure
of $\langle a\rangle$ has index at least $2^{k-1}$ in
$\Z_2^\times$. This observation is not a counterexample to the
proposition, whose congruences are essential and are satisfied in
the submitted construction. The additional requirement $q\mid d+1$
is not needed for this proposition.
\end{remark}

\subsection{A determinant bound that retains finite torsion}

Write $v_p(0)=\infty$, and normalize $|p|_p=p^{-1}$.
If $X$ is generated by $R\ge1$ elements over $\Z_p$, and
$T$ is an endomorphism whose reduction on $X/pX$ is unipotent,
put
\[
 j=1+\min\{k\ge0:p^k\ge R\},\qquad
 b(e,m)=\length_{\Z_p}X/((u^eT)^{p^m}-1)X,
\]
where $v_p(u-1)=c\ge1$, with $c\ge2$ for $p=2$.

\begin{lemma}[Growth in a cyclotomic layer]
\label{ct-growth}
For all $e,m\ge0$,
\[
                    b(e,m)\le b(e,j)+R\max(m-j,0).
\]
\end{lemma}

\begin{proof}
Set $g=u^eT$. Its reduction modulo $p$ is unipotent on a vector
space of dimension at most $R$, so $(g-1)^R X\subset pX$.
If $p^k\ge R$, the binomial theorem in characteristic $p$ gives
$(g^{p^k}-1)X\subset pX$. Raising once more to the $p$th power
gives $(g^{p^{k+1}}-1)X\subset p^2X$: in the binomial expansion,
the first term is divisible by $p^2$, the intermediate terms have
an additional factor $p$, and the last has image in $p^pX$.
This proves the containment at level $j$, and at every larger level.

For any endomorphism $h$ with $(h-1)X\subset p^2X$, put
$Q=1+h+\cdots+h^{p-1}$. Then
$(Q-p)X\subset p^2X$, and hence
$QX\subset pX$ and $pX=QX+p^2X$.
Nakayama's lemma applied to the finitely generated module
$pX/QX$ gives $QX=pX$. As $Q$ commutes with $h-1$,
\[
                    (h^p-1)X=p(h-1)X.
\]
If $A=(h-1)X$, the exact sequence
$0\to A/pA\to X/pA\to X/A\to0$ shows that the cokernel
length increases by at most $R$: $A$ is an image of $X$, so
$\dim A/pA\le R$. The inequality remains true if a length is
infinite. Iteration gives the assertion for $m\ge j$.
For $m<j$, the factorization of $g^{p^j}-1$ by
$g^{p^m}-1$ gives a surjection of the cokernel at level $j$
onto the cokernel at level $m$.
\end{proof}

\begin{proposition}[At most $R$ exceptional weights]
\label{ct-centres}
There are $s_1,\ldots,s_R\in\Z_p$, independent of $e$ and $m$,
such that
\[
 b(e,m)\le
       \sum_{i=1}^R\bigl(c+\max(m,j)+v_p(e-s_i)\bigr).
\]
Consequently, for every integer $Z\ge1$ there is a set
$E_Z\subset\{1,\ldots,Z\}$ of cardinality at most $R$ such that
\[
 b(e,m)\le R\bigl(c+\max(m,j)+\lfloor\log_p Z\rfloor\bigr)
       \quad(1\le e\le Z,\ e\notin E_Z).
\]
In particular every length in this last inequality is finite.
\end{proposition}

\begin{proof}
Choose a surjection $\pi:\Z_p^R\to X$ and lift $T^{p^j}$ to a
matrix $B\in M_R(\Z_p)$, so that $\pi B=T^{p^j}\pi$.
For $a\in\Z_p$, the cokernel of $aB-1$ surjects onto that of
$aT^{p^j}-1$. Smith normal form gives
\[
 \length X/(aT^{p^j}-1)X\le v_p\det(aB-1).
\]
The convention $v_p(0)=\infty$ handles a singular matrix.
No torsion subgroup of $X$ has been discarded.

In a finite splitting field over $\Q_p$, let $z_i$ be the
eigenvalues of $B$, counted with multiplicity. The elementary
lifting-the-exponent formula gives
\[
 |u^{p^j a}-u^{p^j b}|_p
          =p^{-(c+j)}|a-b|_p\qquad(a,b\in\Z_p).
\]
For completeness, for integral $a-b$ prime to $p$ the first term
of the binomial expansion has strictly smallest valuation; raising
to the $p$th power increases that valuation by exactly one.
The condition $c\ge2$ when $p=2$ ensures strictness there.
Approximation by integers gives the displayed identity on $\Z_p$,
and defines the continuous powers used in it.

For each $z_i$ choose $s_i\in\Z_p$ that minimizes
$|z_i-u^{-p^j s}|_p$. Such a point exists by compactness.
The ultrametric inequality and the minimizing property give
\[
 |z_i-u^{-p^j e}|_p
       \ge |u^{-p^j e}-u^{-p^j s_i}|_p
       =p^{-(c+j)}|e-s_i|_p.
\]
Since $|u|_p=1$, the left side equals
$|u^{ep^j}z_i-1|_p$. Multiply these inequalities and use
$\det(u^{ep^j}B-1)=\prod_i(u^{ep^j}z_i-1)$.
This proves the bound at level $j$; Lemma~\ref{ct-growth}
proves it at every level.

Put $k=\lfloor\log_p Z\rfloor$. Each ball
$\{e\in\{1,\ldots,Z\}:|e-s_i|_p<p^{-k}\}$ contains at most
one integer. Two distinct integers in it would have difference
divisible by $p^{k+1}$, although its ordinary absolute value is
less than $p^{k+1}$. The union of these balls is the required
$E_Z$. Outside it, each $v_p(e-s_i)\le k$.
\end{proof}

For the module in Proposition~\ref{ct-descent}, the unipotence
hypothesis is automatic once it is finitely generated: the action
on the finite space $X_p/pX_p$ is continuous and factors through
a finite quotient of $\Gamma\simeq\Z_p$, hence through a
$p$-group. If $T^{p^a}=1$ on that space, then
$(T-1)^{p^a}=0$. The continuity follows either from the compact
dual construction or directly from the finite orbit of each
continuous finite-valued character. Thus all hypotheses in the
two preceding results hold for the actual cyclotomic modules.

\begin{example}[Why two tempting replacements fail]
Let $X=\Z/p^a\Z$, $a\ge1$, $T=1$, and $u=1+p^c$, with
the same restrictions on $c$. Then
\[
       b(e,m)=\min\{a,c+m+v_p(e)\}\quad(e\ge1).
\]
The equality follows from the lifting-the-exponent formula.
Passing to $X$ modulo its torsion would give the zero module and
would incorrectly make this positive length vanish. Also, the
strict inequality in the exceptional balls is necessary:
with $s=1$, $k\ge1$, and $Z=1+p^k<p^{k+1}$, both $1$ and $1+p^k$
have distance at most $p^{-k}$ from $s$. These examples refute
the proposed simplifications, not the actual results above or the
qualified source statements.
\end{example}

\subsection{Avoiding all the exceptional weights on one interval}

\begin{lemma}[A common multiplicative interval]
\label{ct-common-interval}
Let $E$ be a finite set of positive integers with $\#E\le R$,
let $C_t\ge\eta>0$, let $b>1$, and put
\[
             J=b^{R+1}(C_t/\eta)^R.
\]
For every $X>1$, there is $Y$ satisfying
\[
                 X^{1/J}\le Y,\qquad Y^b\le X,
\]
such that, simultaneously for every $x\in[Y,Y^b]$ and $e\in E$,
\[
                     e\le x^{C_t}\ \Longrightarrow\ e\le x^\eta.
\]
\end{lemma}

\begin{proof}
We prove the corresponding real assertion. Given $U>0$, $B\ge1$,
and a set $A\subset\R$ of at most $R$ points, there is
$y\ge U/[b(bB)^R]$ with $by\le U$ such that
$a\le v\Longrightarrow Ba\le v$ for all $a\in A$ and
$v\in[y,by]$.
Induct on $R$. If no point lies in $(U/(bB),U]$, take $y=U/b$:
every relevant point satisfies $Ba\le U/b$.
Otherwise choose a point $a$ in that interval, remove it, and
apply the induction hypothesis with $U'=U/(bB)$.
The resulting interval lies below $a$, so the removed point
imposes no condition; the bound on $y$ is exactly the required one.
For the empty set the choice $y=U/b$ starts the induction.
Now set $U=\log X$, $B=C_t/\eta$,
$A=\{\log(e)/C_t:e\in E\}$, and $Y=e^y$.
Taking logarithms gives the stated implications and endpoints.
\end{proof}

For $b=2$, this gives $J=2^{R+1}(C_t/\eta)^R$, improving the
factor $2^{R+2}$ in the corresponding source argument. More
significantly, the exceptional sets for all coefficient primes are
united \emph{before} the lemma is applied. Intersecting intervals
chosen separately would not prove that a common interval exists.

\subsection{The estimate at the exceptional weights}

We state the coarse estimate used at small weights as an explicit
hypothesis. There is an absolute $C_K\ge1$ such that, for every
totally complex number field $F$ containing $\mu_p$ (and $\mu_4$
when $p=2$) and every $e\ge1$, the group
$H^2(G_{F,S},\Z_p(e+1))$ is finite and
\begin{equation}
 \label{ct-coarse-hypothesis}
 \log\#H^2(G_{F,S},\Z_p(e+1))
 \le \exp\bigl((2e[F:\Q])^{C_K}\log(2|D_F|)\bigr).
\end{equation}
Here $D_F$ is the absolute field discriminant. Poitou--Tate
duality identifies the dual of $\mathcal S_p(F,e)$ with the
subgroup of this $H^2$ group whose localizations vanish. Thus
\begin{equation}
 \label{ct-small-weight}
 \lambda_p(F,e)\log p
 \le \log\#H^2(G_{F,S},\Z_p(e+1)).
\end{equation}
We use this precise consequence of duality; it is a statement
about a subgroup and includes finiteness. Theorem~\ref{kt-fourteen}
proves \eqref{ct-coarse-hypothesis} with $C_K=14$, using the
stated classical comparison theorems. Finite generation alone
would not give this order estimate.

\begin{proposition}[The common interval under the coarse estimate]
\label{ct-conditional-interval}
Assume \eqref{ct-coarse-hypothesis}. Fix $M$, a finite set $P$ of
odd primes, and for each $p\in P$ a fixed abelian field $K_p$
containing $\mu_p$. Fix $C_t\ge1$, $A>0$, and $0<\epsilon<1$.
Choose generator bounds $R_2,R_p\ge1$ as in
Proposition~\ref{ct-uniform-generators}; set
\[
 R_*=R_2+\sum_{p\in P}R_p,\quad
 \eta=\frac{\epsilon}{2C_K},\quad
 J=2^{R_*+1}(C_t/\eta)^{R_*}.
\]
For each coefficient prime use the fixed layer
$j_p=1+\min\{k:p^k\ge R_p\}$.
For all sufficiently large $X$, uniformly for the primes $q,d$
in Proposition~\ref{ct-uniform-generators} satisfying
\[
                     \log(qd)\le X^{\epsilon/(4J)},
\]
there is $Y$ with $X^{1/J}\le Y$ and $Y^2\le X$ such that,
for all $x\in[Y,Y^2]$, all integers $0\le h\le A(1+\log x)$,
and all integers $1\le t\le x^{C_t}$,
\begin{equation}
 \label{ct-summed-bound}
 (\log2)\sum_{1\le e<t}\lambda_2(E_h,e)
   +\sum_{p\in P}(\log p)
          \sum_{1\le e<t}\lambda_p((K_p)_{j_p},e)
        \le \exp(x^{4\epsilon/5}),
\end{equation}
where
$E_h=M\Q(\mu_{2^{h+2}})^+(\sqrt{-q},\sqrt{-d})$.
All lengths in this inequality are finite. The threshold for $X$
may depend on the fixed data and on $\epsilon$, but not on $q,d$.
\end{proposition}

\begin{proof}
Put $L=M(i,\sqrt{-q},\sqrt{-d})$. Its degree is at most
$8d_M$, so the largest $c$ with $\mu_{2^c}\subset L$ is bounded
in terms of $M$: $2^{c-1}\le[L:\Q]\le8d_M$.
Its fixed layer $L_{j_2}$ has bounded degree, and
\begin{equation}
 \label{ct-discriminant}
        \log(2|D_{L_{j_2}}|)\ll_M 1+\log(qd),
\end{equation}
where the constant may also depend on the fixed $j_2$.
Here is a direct verification of this last uniformity. Choose a
fixed integer $c_0$ bounding $c$, and put
$B=M(i,\mu_{2^{c_0+j_2}})$. The field $L_{j_2}$ is a subfield
of $B(\sqrt{-q},\sqrt{-d})$. For a quadratic extension
$F(\sqrt a)/F$ with integral $a$, the relative discriminant
divides $(4a)$: the order $\OO_F[\sqrt a]$ has that
discriminant and passage to the maximal order divides it by the
square of an index. If the extension is trivial the same upper
bound holds. Applying the discriminant tower formula twice gives
\[
 |D_{B(\sqrt{-q},\sqrt{-d})}|
 \le |D_B|^4(4q)^{2[B:\Q]}(4d)^{2[B:\Q]}.
\]
The discriminant tower formula also bounds the discriminant of
any subfield by that of this field. This proves
\eqref{ct-discriminant}. The odd-prime fixed layers have fixed
degree and discriminant.

Apply Proposition~\ref{ct-centres} with
$Z=\lfloor X^{C_t}\rfloor$ to all these modules. Unite their
exceptional sets; the union has at most $R_*$ elements. Lemma
\ref{ct-common-interval} gives the required $[Y,Y^2]$. For any
$x$ in it, every exceptional weight under consideration is at
most $x^\eta$. For a nonexceptional weight the length is
$O(1+\log X+h)=O(1+\log x)$, since
$\log X\le J\log x$. The constants are uniform in $q,d$.

For an exceptional weight, the growth lemma reduces to the fixed
layer, at a cost $O(1+\log x)$. At that layer, use
\eqref{ct-small-weight} and \eqref{ct-coarse-hypothesis}.
As $e\le x^\eta$, the degree is bounded, and
$X^{1/J}\le x$, equation \eqref{ct-discriminant} gives
\[
 (2e[F:\Q])^{C_K}\log(2|D_F|)
       \le C x^{\epsilon/2}x^{\epsilon/4}
       =C x^{3\epsilon/4}.
\]
Thus each exceptional length, multiplied by $\log p$, is at most
$C(1+\log x)+\exp(Cx^{3\epsilon/4})$.

To pass from the $2$-power layers of $L$ to $E_h$, observe that
$E_h(i)=L(\mu_{2^{h+2}})=L_m$ for some $m\le h+2$ and
$[E_h(i):E_h]\le2$. Restriction carries the strict local
conditions to the same strict conditions. Its kernel is a
subgroup of
$H^1(\operatorname{Gal}(E_h(i)/E_h),A_e^{G_{E_h(i)}})$,
which has order at most $2$. Indeed, this $H^1$ is killed by $2$
and is a subquotient of a subgroup of $\Q_2/\Z_2$, so its
order is at most $2$. Hence
\[
                    \lambda_2(E_h,e)\le1+\lambda_2(L_m,e).
\]

In each prime-wise sum use that prime's own exceptional set.
There are at most $R_*$ exceptional terms across these sums,
counted with multiplicity, and at most
$O(x^{C_t})$ other terms. The left side of
\eqref{ct-summed-bound} is therefore bounded by
\[
 Cx^{C_t}(1+\log x)+C\exp(Cx^{3\epsilon/4}).
\]
For fixed $C,C_t,\epsilon>0$ this is at most
$\exp(x^{4\epsilon/5})$ once $x$ is sufficiently large:
$3/4<4/5$ and every fixed multiple of $\log x$ is
$o(x^{4\epsilon/5})$. Increasing the threshold for $X$ makes
this hold throughout the chosen interval. The determinant bounds
and the coarse estimate separately give finiteness in the two
cases, so no infinite length has been silently converted to zero.
\end{proof}

This proves the common-interval mechanism and its uniformity in
the auxiliary quadratic fields. The initially displayed coarse
hypothesis is supplied by Theorem~\ref{kt-fourteen}.
The full cohomological estimate needed
for GRH additionally controls all the other primes and all relevant
character components. Those other estimates remain to be reconstructed. In
particular, this section does not replace the automorphy and
deformation arguments used for those components.


